% sections/02-technical-requirements.tex — Technical Requirements

\section{Technical Requirements}

\begin{constraintbox}[HARD CONSTRAINTS --- NON-NEGOTIABLE]
\begin{itemize}
  \item Makefile or CMake with \make, \maketest, \makeclean{} targets
  \item \cppsev{} or later (\code{g++} / \code{clang++})
  \item Compiles without errors on submission
  \item \code{clang-format} configuration file included
\end{itemize}
\vspace{4pt}
\centering\textbf{Missing any of these = cannot grade your project.}
\end{constraintbox}

% ── Programming Language ──
\subsection{Programming Language}

\begin{itemize}
  \item \textbf{Primary Language:} \cpp{} (\cppsev{} or later)
  \item \textbf{Compiler:} \code{g++} or \code{clang++} with \code{-std=c++17}
  \item \textbf{Platform:} Linux (Ubuntu 20.04+) or Windows
\end{itemize}

% ── Build System ──
\subsection{Build System}

Your project must include a \filename{Makefile} or \filename{CMakeLists.txt}. The following build targets are required:

\begin{center}
\begin{tabular}{l l}
\hline
\rowcolor{tableheader} \textbf{Command} & \textbf{Description} \\
\hline
\make{} / \code{cmake --build .} & Build the project \\
\makerun & Build and run the application \\
\maketest & Run all tests \\
\makedeps & Install dependencies \\
\makeclean & Remove build artifacts \\
\makedocs & Generate documentation \\
\hline
\end{tabular}
\end{center}

% ── Code Formatting ──
\subsection{Code Formatting}

All source code must be formatted using \code{clang-format}. Include a \filename{.clang-format} configuration file in your project root.

Recommended styles:
\begin{itemize}
  \item Google \cpp{} Style
  \item LLVM Style
  \item Mozilla Style
\end{itemize}

To format all source files:

\begin{lstlisting}[style=bash]
clang-format -i src/*.cpp include/*.h
\end{lstlisting}

% ── Documentation Style ──
\subsection{Documentation Style}

All code documentation must follow Doxygen format:

\begin{lstlisting}[style=cpp]
/**
 * @brief Inserts a value into the binary search tree.
 *
 * @param value The value to insert.
 * @return true if the value was inserted successfully.
 * @throws std::invalid_argument if the value already exists.
 * @pre The tree is in a valid state.
 * @post The tree contains the new value and remains balanced.
 */
bool insert(int value);
\end{lstlisting}

Include a \filename{Doxyfile} in your \filepath{docs/} directory. Generate documentation with:

\begin{lstlisting}[style=bash]
doxygen docs/Doxyfile
\end{lstlisting}
